\section*{Chapitre \ref{ch:in:bank-quant} --- Quant de banque}

\begin{solution}{exo:in:bank-quant:1}
De 128\,300~\$ à 158\,100~\$~: 23.2~\% en dollars. Les prix ont augmenté de $135.8312/114.325=1.188$, donc la hausse réelle est
$1.232/1.188-1=3.7~\%$.
\end{solution}

\begin{solution}{exo:in:bank-quant:2}
Organisation (a)~: les deux unités rendent compte à des membres différents de la direction générale (le directeur des risques et le
responsable des marchés), et le validateur ne dépend ni du développeur ni du propriétaire.
\end{solution}

\begin{solution}{exo:in:bank-quant:3}
Niveau un~: une validation complète chaque année (400 heures), pas de revue séparée, et 10~\% de 400 pour les changements importants~: 440 heures.
Niveau trois~: $60/3+8\times\frac23+0.1\times60=31.3$ heures.
\end{solution}

\begin{solution}{exo:in:bank-quant:4}
Avec un cycle d'un an pour le niveau deux, chaque modèle de ce niveau coûte $160+16=176$ heures au lieu de 104~: l'inventaire demande 82.1
validateurs au lieu de 63.9.
\end{solution}

\begin{solution}{exo:in:bank-quant:5}
Le terme lié aux changements double, passant de 13\,498 à 26\,996 heures~: 72.4 validateurs.
\end{solution}

\begin{solution}{exo:in:bank-quant:6}
Le stratégiste a intérêt à ce que le modèle soit approuvé, connaît ce pour quoi il a été conçu plutôt que ce qu'il ne fait pas,
et réviserait des hypothèses qu'il a lui-même choisies. Chaque règle du chapitre (la notion d'«~intérêt~» des recommandations américaines de 2011, la déclaration britannique sur les
lignes hiérarchiques séparées, les dispositions de la BCE) l'excluent~; les connaissances du stratégiste entrent sous forme de documentation et
de tests que le validateur examine.
\end{solution}

\begin{solution}{exo:in:bank-quant:7}
Chaque tirage répartit en niveaux 1\,000 triplets de scores aléatoires et applique la méthode~: la moyenne est de 64.0 validateurs et l'intervalle entre le 5\up{e}~centile
et le 95\up{e}~centile va de 60.1 à 67.9. L'aléa de l'inventaire déplace la réponse d'environ trois validateurs dans un sens ou dans l'autre~;
le choix des heures la déplace bien davantage.
\end{solution}

\begin{solution}{exo:in:bank-quant:8}
Les validateurs rendent compte au responsable de l'analytique quantitative, qui possède la bibliothèque qu'ils examinent~: ils dépendent du
développeur, donc \texttt{independence} renvoie \texttt{fail}, quel que soit le nom de l'équipe. L'organisation (b) suppose deux unités
qui ne se rejoignent qu'à un membre de la direction générale.
\end{solution}

\begin{solution}{pb:in:bank-quant:1}
\begin{enumerate}
\item Comme dans les quatre définitions du chapitre~: le constructeur du modèle et de l'outil du desk~; le constructeur de la bibliothèque de valorisation partagée~;
le constructeur des modèles de risque et de capital~; le relecteur indépendant des modèles des autres.
\item Goldman Sachs~: les \omterm{def:in:bank-quant:strat}{stratégistes de desk} (desk strats) «~siègent dans les salles de marché et apportent de la valeur aux clients internes \& externes grâce
à des modèles de pointe, à l'analyse prédictive et à des outils de trading de haute qualité~».
\item Pour qu'une même transaction ait un seul prix dans le système de trading, le moteur de risque et la comptabilité.
\item De nombreux modèles de risque et de capital requièrent l'approbation d'un superviseur~; un changement important la requiert à nouveau.
\item Entre le desk et le risque~: il simule les expositions comme un \omterm{def:in:bank-quant:risk}{quant de risque}, valorise les charges comme un stratégiste, et ses
chiffres entrent dans la comptabilité.
\item Les validateurs ne doivent avoir aucun intérêt dans le résultat~; les lignes hiérarchiques peuvent favoriser l'indépendance, mais elle s'apprécie
à partir des actions et des résultats~; chaque modèle est revu au moins une fois par an.
\item La révision 2026 des recommandations interagences (SR~26-2~; OCC Bulletin 2026-13)~: fondée sur les risques, le calendrier et la fréquence de
validation variant selon le modèle.
\item Des lignes hiérarchiques séparées pour les validateurs, d'une part, et pour les développeurs et les propriétaires, et une hiérarchisation selon l'importance et la complexité.
\item (a) Des membres différents de la direction générale~; (b) le même membre, des unités différentes~; (c) un personnel distinct au sein d'une même unité.
\item L'audit interne évalue si la gestion du risque de modèle est rigoureuse et efficace~; il ne duplique pas le développement
ni la validation.
\item $H=\sum_t n_t(F_t/k_t+R_t(1-1/k_t)+cF_t)$ heures~; $H/h$ validateurs.
\item Deux fois le score d'importance, plus la complexité et l'incertitude~: 10 ou plus correspond au niveau un (annuel), 7 à 9 au niveau deux
(tous les deux ans), en dessous au niveau trois (tous les trois ans).
\item 102, 403 et 495.
\item 43.9~\%, pour 10.2~\% des modèles~: ils sont revalidés en totalité chaque année, aux heures les plus élevées.
\item 18.1 validateurs (82.1 contre 63.9).
\item 63.9 validateurs~; environ 2.6 de plus pour chaque point de pourcentage de modèles de niveau un.
\item À titre illustratif~: les probabilités des scores, les heures, le taux de changement et les heures productives. Les règles proviennent de la revue annuelle
prévue par les recommandations américaines de 2011, de la validation annuelle des modèles internes du guide de la BCE et de la hiérarchisation de la déclaration britannique~; les
cycles par niveau relèvent de la politique illustrative du livre~6.
\item Organisation (a), la plus robuste.
\item L'audit interne, qui vérifie que la fonction fonctionne~; et les superviseurs, qui examinent le cadre de la banque.
\item Avec la hiérarchisation actuelle de la banque, l'inventaire demande environ 64 validateurs, et chaque point de pourcentage de modèles basculés
au niveau un en ajoute environ trois. L'équipe rend compte au directeur des risques, séparément des développeurs et des
propriétaires, ce qui respecte chacune des règles citées dans le chapitre.
\end{enumerate}
\end{solution}

\begin{solution}{iq:in:bank-quant:1}
Comparer avec une implémentation indépendante sur une grille (strikes, maturités, volatilités, taux), vérifier les cas limites (volatilité nulle,
maturité nulle, options profondément dans ou hors de la monnaie), la parité call-put, les grecques comparées aux différences finies, et
le comportement pour des entrées extrêmes.
\iqlookfor{un indice de référence indépendant, des cas limites et des identités, et non un seul prix de test.}
\end{solution}

\begin{solution}{iq:in:bank-quant:2}
Chercher d'abord à comprendre pourquoi avant de modifier quoi que ce soit~: les mêmes entrées (surface de volatilité, suivi des barrières, rebate, dividendes)~? Le
modèle (volatilité locale ou stochastique, qui ne valorisent pas les barrières de la même façon)~? Un écart avec le prix du marché est un indice
sur le modèle à comprendre, et non un prix à recopier.
\iqlookfor{séparer les entrées du risque de modèle.}
\end{solution}

\begin{solution}{iq:in:bank-quant:3}
Deux et demi sont attendus~; sept ou plus a une probabilité de 1.4~\% sous un modèle correct, ce qui constitue un indice contre lui, et
il tombe dans la zone jaune de Bâle (cinq à neuf). Vérifier d'abord les erreurs de données et de positions, ainsi que le partitionnement (clustering) des dépassements dans une
période de stress, puis le modèle.
\iqlookfor{un test binomial et un examen des causes.}
\end{solution}

\begin{solution}{iq:in:bank-quant:4}
Mesurer le changement sur les positions de chaque desk avant la mise en production, l'expliquer, le faire valider comme un changement de modèle, convenir d'une date
avec les desks et le product control, et mettre en production en gardant l'ancienne version disponible pour comparaison.
\iqlookfor{analyse d'impact et contrôle des changements.}
\end{solution}

\begin{solution}{iq:in:bank-quant:5}
Son usage prévu et les décisions qu'il alimente~; ses limites telles qu'énoncées par le développeur~; les données et les hypothèses~; et les tests réalisés,
et l'indice de référence auquel ils ont été comparés. Un document qui ne précise pas à quoi le modèle ne doit pas servir constitue une constatation.
\iqlookfor{usage, limites et éléments de preuve.}
\end{solution}

\begin{solution}{iq:in:bank-quant:6}
L'écart relève d'un risque de modèle que la surface actuelle ne peut pas résoudre, car un cliquet dépend de la volatilité forward.
Valoriser avec l'un des deux, réserver l'écart ou le présenter comme une incertitude de modèle, essayer un modèle à volatilité stochastique, et fournir
aux validateurs les deux calibrations et l'écart.
\iqlookfor{dépendance à la volatilité forward et réserve de modèle.}
\end{solution}
